\section{Decision Audit Framework}
\label{sec:framework}

We compare two eligibility rules applied after forecast evaluation. The rules do not retrain models or create a new ranking objective; they audit how a stated requirement changes the set of eligible cells.

\subsection{Error-based eligibility: Rule A}
For each station--model--horizon cell,
\begin{equation}
\textbf{Rule A:}\qquad \mathrm{Skill}>0 \quad\mathrm{and}\quad \mathrm{DM}_{BH}\ \text{significant}.
\end{equation}
Skill is persistence-relative RMSE skill. The significance flag is based on the two-sided Harvey--Leybourne--Newbold modified Diebold--Mariano comparison with Benjamini--Hochberg adjustment, as implemented in \texttt{scripts/05\_dm\_significance.py} \cite{diebold1995,harvey1997}.

\subsection{Fidelity-aware eligibility: Rule B}
\begin{equation}
\textbf{Rule B:}\qquad \text{Rule A}\quad\mathrm{and}\quad \alpha\geq0.50\quad\mathrm{and}\quad \mathrm{Recall}_{P75}\geq0.20.
\end{equation}
The two thresholds are operational criteria chosen for this audit. They are not estimated optima, universal standards, or claims about acceptable performance in other applications. A decision change is a cell that satisfies Rule A but not Rule B.

We also define a descriptive discordant case as a Rule-A cell with $\mathrm{Recall}_{P75}\geq0.20$ and $\alpha<0.50$. This isolates cells that pass the error and episode-sensitivity checks while failing the variance-retention check.
